\section{Context Type System}
\label{sec:context-typing}
% In this section, we introduce the core language and the context type system for it.

% \subsection{Core Language}

\begin{figure}[h]
  {\small
    \begin{alignat*}{2}
      \text{\textbf{Variables }}& \quad &\quad& x, y, z, \vnu, ...
      \\\text{\textbf{Data Constructors }}& \quad &d ::= \quad & () ~|~ \Code{true} ~|~ \Code{false} ~|~ \Code{O} ~|~ \Code{S} ~|~ \Code{Cons} ~|~ \Code{Nil} ~|~ ...
      \\\text{\textbf{Operations}}& \quad & \primop ::= \quad & d ~|~ {+} ~|~ {-} ~|~ {==} ~|~ {<} ~|~ {\leq} ~|~ ...
      \\ \text{\textbf{Constants }}& \quad &c ::= \quad & () ~|~ \mathbb{B} ~|~ \mathbb{Z} ~|~ d\ \overline{c}\ ...
      \\ \text{\textbf{Values}}& \quad  & v ::= \quad & c ~|~ x ~|~ \zlam{x}{s}{e} ~|~ \zfix{f}{s}{x}{s}{e}
      \\ \text{\textbf{Expressions}}& \quad & e ::=\quad & v ~|~ \zlet{x{:}s}{e}{e} ~|~ \primop\ \overline{v} ~|~  v\;v ~|~ \match{v} \overline{d\ \overline{y} \to e}
      \\\text{\textbf{Base Types}}& \quad & b  ::= \quad &  \Code{unit} ~|~ \Code{bool} ~|~ \Code{nat} ~|~ \List[b] ~|~ ...
      \\\text{\textbf{Basic Types}}& \quad & s  ::= \quad &  b ~|~ s\sarr s
      \\\text{\textbf{Basic Type Context}}& \quad & \Sigma ::= \quad &  \emptyset ~|~ x{:}s ~|~ \Sigma, \Sigma
    \end{alignat*}}
  \caption{Syntax of \langname{} \BD{Why do we include both constants
      and data constructors?}}
  \label{fig:core-lang-syntax}
  \end{figure}

  % \paragraph{Terms, basic types, and basic type contexts}
  We formalize our proposed context type system for the core language
  (\langname{}) shown in \autoref{fig:core-lang-syntax}.  This
  language is a standard call-by-value lambda calculus with inductive
  datatypes, pattern matching (the following uses conditional
  expressions as syntactic sugar for \Code{match}), and recursive
  functions. \BD{Are we assuming that built-in operations are
    deterministic? Do we even need these?}  Expressions are written in
  monadic normal form (MNF)~\cite{JO94}, a variant of A-Normal Form
  (ANF)~\cite{FSDF93} that allows nested let-bindings.  The language's
  basic types (i.e., $s$) consist of base types ($b$) and function
  types (i.e., $s\sarr s$). Both \textbf{let}-expressions and function
  abstractions explicitly annotate bound variables with a basic type
  (e.g., $\zlam{x}{s}{e}$); we omit these annotations for brevity,
  when appropriate.  A basic type context $\Sigma$ maps variables to
  their basic types; we assume that all bound variables in $\Sigma$
  are distinct.

\begin{figure}[h]
  { \small
    \begin{alignat*}{2}
      \text{\textbf{Qualifiers}}& \quad & \phi ::= \quad & c ~|~ x ~|~ \primop\;\overline{v} ~|~ \bot ~|~ \top  ~|~ \neg \phi ~|~ \phi \land \phi ~|~ \phi \lor \phi ~|~  \phi \impl \phi ~|~ \forall x{:}b. \phi ~|~ \exists x{:}b. \phi
      \\ \text{\textbf{Refinement Types}}& \quad & \tau ::= \quad &
      \ort{b}{\phi} ~|~ \urt{b}{\phi} ~|~ \tau \ctinter \tau ~|~ \tau
      \ctoplus \tau ~|~ x{:}\tau \ctsarr \tau ~|~ x{:}\tau \ctwand
      \tau %~|~ \tau \ctunion \tau
      % ~|~ \ctpers \tau
      \\ \text{\textbf{Type Contexts}}& \quad & \Gamma ::= \quad & \emptyset ~|~ x{:}\tau ~|~ \Gamma \ctstar \Gamma ~|~ \Gamma \ctcomma \Gamma ~|~ \Gamma \ctoplus \Gamma
      % \\ \text{\textbf{Contexts With Holes}}& \quad & \Delta ::= \quad & \Box ~|~ \Gamma ~|~ \Delta \ctstar \Delta ~|~ \Delta \ctcomma \Delta ~|~ \Delta \ctoplus \Delta
    \end{alignat*}}
  %% \vspace*{-.1in}
  %% {\small
  %% \begin{align*}
  %%   \text{\textbf{Type erasure}} \quad & \eraserf{\ort{b}{\phi}} \doteq b \quad \eraserf{\urt{b}{\phi}} \doteq b \quad \eraserf{\tau_1 \ctinter \tau_2} \doteq \eraserf{\tau_1} \quad \eraserf{\tau_1 \ctunion \tau_2} \doteq \eraserf{\tau_1}
  %%   \\&\eraserf{\tau_1 \ctoplus \tau_2} \doteq \eraserf{\tau_1} \quad  \eraserf{x{:}\tau_x \ctsarr \tau} \doteq \eraserf{\tau_x} \sarr \eraserf{\tau} \quad
  %%   \\& \eraserf{\emptyset} \doteq \emptyset \quad \eraserf{x{:}\tau} \doteq x{:}\eraserf{\tau} \quad \eraserf{\Gamma_1 \ctcomma \Gamma_2} \doteq \eraserf{\Gamma_1}, \eraserf{\Gamma_2} \quad \eraserf{\Gamma_1 \ctstar \Gamma_2} \doteq \eraserf{\Gamma_1}, \eraserf{\Gamma_2}  \quad \eraserf{\Gamma_1 \ctoplus \Gamma_2} \doteq \eraserf{\Gamma_1}
    % \\\text{\textbf{Type lifting}} \quad & \liftrf{b} \doteq \ort{b}{\phi} \quad \liftrf{s_x \sarr s} \doteq x{:}\liftrf{s_x} \ctsarr \liftrf{s}
    % \\&\liftrf{\emptyset} \doteq \emptyset \quad \liftrf{x{:}s} \doteq x{:}\liftrf{s} \quad \liftrf{\Gamma_1 , \Gamma_2} \doteq \liftrf{\Gamma_1} \ctcomma \liftrf{\Gamma_2}
%  \end{align*}}
  \caption{Context Type Syntax}
  \label{fig:context-types-syntax}
\end{figure}

The syntax of context types is shown in
\autoref{fig:context-types-syntax}. As in other refinement type
systems, \emph{context types} refine basic types with qualifiers
$\phi$; but each refinement is ascribed either a demonic and angelic
modality ($\otparenth{...}$ and $\utparenth{...}$, resp.). Context
types additionally include bunched implication ($\ctwand$), disjoint
sums ($\ctoplus$), and standard type intersection ($\ctinter$)
operators. Our type system shares the same intuition underlying
bunched implication and refinement types, generally; however, its
technical development differs significantly from both due to its
support for angelic contexts and mixed demonic/angelic reasoning.

\begin{figure}[h!]
  {\small
  {\normalsize \begin{flalign*}
    &\text{\textbf{Well-formedness rules}} && \fbox{$\Sigma \wfvdash \tau \quad \Sigma \wfvdash \Gamma$}
  \end{flalign*}}
% \mprooftr{28mm}
% {$\Gamma \basicvdash \tau$ \quad $\wfvdash \Gamma$}
% {WfType}
% {$\Gamma \wfvdash \tau$}
% \quad
\mprooftr{35mm}
{$\Sigma \wfvdash \Gamma_1$ \quad $\Sigma, \eraserf{\Gamma_1} \wfvdash \Gamma_2$}
{WfComma}
{$\Sigma \wfvdash \Gamma_1 \ctcomma \Gamma_2$}
\quad
\mprooftr{35mm}
{$\Sigma \wfvdash \Gamma_1$ \quad $\Sigma \wfvdash \Gamma_2$}
{WfStar}
{$\Sigma \wfvdash \Gamma_1 \ctstar \Gamma_2$}
% \mprooftr{53mm}
% {$\forall r. r \models \denot{\Gamma} \impl \forall e. r \models \denot{\tau_1}\;e \impl r \models \denot{\tau_2}\;e$}
% {Sub}
% {$\Gamma \vdash \tau_1 <: \tau_2$}
% \quad
% \mprooftr{50mm}
% {$\forall r. \denot{\Gamma_1}\;r \impl \exists r'. \restrictTo{r}{X} \sqsubseteq r' \land \denot{\Gamma_2}\;r'$ }
% {CtxSub}
% {$\Gamma_1 \ctsubseteq{X} \Gamma_2$}
	  }
  \caption{Selected auxiliary well-formedness relations\footnote{Well-formedness rules for types
      and contexts check that qualifiers are well-typed with respect to the basic
      type structure.  The $\eraserf{.}$ operator converts context types $\tau$ to basic types $s$
      by stripping away all refinement information.  Details are provided in the supplemental material.}}
  \label{fig:auxiliary-rules}
  \vspace*{-.1in}
\end{figure}

\paragraph{Well-formedness} Similar to other refinement type systems,
our type system depends on the well-formedness relations shown in
\autoref{fig:auxiliary-rules}. The well-formedness of types
($\Sigma \wfvdash \tau$) and type contexts ($\Sigma \wfvdash \Gamma$)
mostly requires type qualifiers to be well-typed in the basic type
system and binding variables in type contexts to be unique. We would
like to highlight that rule \textsc{WfComma} allows the second context
to refer to variables in the first context, while rule \textsc{WfStar}
requires that $\Gamma_1$ and $\Gamma_2$ are completely
independent. For example, within a type context
$\Gamma_0 \ctcomma (\Gamma_1 \ctstar \Gamma_2)$, context $\Gamma_0$ is
closed, and $\Gamma_1$ and $\Gamma_2$ can only refer to variables in
$\Gamma_0$.

\begin{figure}[h!]
  {\small
  {\normalsize \begin{flalign*}
    &\text{\textbf{Subtyping rules}} && \fbox{$\Gamma \vdash \tau \subty \tau$ \quad $\Gamma \ctsubseteq{X} \Gamma$ }
    \\&\text{\textbf{Typing rules}} && \fbox{$\Gamma \vdash e : \tau$}
  \end{flalign*}}
  \mprooftr{40mm}
  {$\Gamma \vdash e_1 : \tau_1 $ \quad $\Gamma \ctcomma x{:}\tau_1 \vdash e_2 : \tau_2$}
  {T-Let}
  {$\Gamma \vdash \zlet{x}{e_1}{e_2} : \tau_2$}
  \quad
  \mprooftr{40mm}
  {$\Gamma_1 \vdash e_1 : \tau_1 $ \quad $\Gamma_2 \ctstar x{:}\tau_1 \vdash e_2 : \tau_2$}
  {T-LetD}
  {$\Gamma_1 \ctstar \Gamma_2 \vdash \zlet{x}{e_1}{e_2} : \tau_2$}
  \\[.8em]
  \mprooftr{42mm}
  {$\eraserf{\tau_x} = s_x$ \quad $\Gamma \ctcomma x{:}\tau_x \vdash e : \tau$}
  {T-Lam}
  {$\Gamma \vdash \zlam{x}{s_x}{e} : x{:}\tau_x \ctsarr \tau$}
  \quad
  \mprooftr{42mm}
  {$\eraserf{\tau_x} = s_x$ \quad  $\Gamma \ctstar x{:}\tau_x \vdash e : \tau$}
  {T-LamD}
  {$\Gamma \vdash \zlam{x}{s_x}{e} : x{:}\tau_x \ctwand \tau$}
  \\[.8em]
  \mprooftr{40mm}
  {$\Gamma \vdash v_1 : x{:}\tau_x \ctsarr \tau$ \quad
  $\Gamma \vdash v_2 : \tau_x$ }
  {T-AppFun}
  {$\Gamma \vdash v_1\;v_2 : \tau[x \mapsto v_2]$}
  \quad
  \mprooftr{46mm}
  {$\Gamma_1 \vdash v_1 : x{:}\tau_x \ctwand \tau$ \quad
  $\Gamma_2 \vdash v_2 : \tau_x$  }
  {T-AppFunD}
  {$\Gamma_1 \ctstar \Gamma_2 \vdash v_1\;v_2 : \tau[x \mapsto v_2]$}
  \\[.8em]
  \mprooftr{30mm}
  {$\Gamma \vdash e : \tau_1$ \quad $\Gamma \vdash \tau_1 \subty \tau_2$}
  {T-Sub}
  {$\Gamma \vdash e : \tau_2$}
  \quad
  \mprooftr{40mm}
  {$\Gamma_2 \vdash e : \tau$ \quad $\Gamma_1 \ctsubseteq{\Code{FV}(e) \cup \Code{FV}(\tau)} \Gamma_2$}
  {T-CtxSub}
  {$\Gamma_1 \vdash e : \tau$}
  \quad
  \mprooftr{14mm}
  {$ $}
  {T-Var}
  {$x{:}\tau \vdash x : \tau$}
  \\[.8em]
  \mprooftr{43mm}
  {$ \emptyset \basicvdash c : b $}
  {T-Const}
  {$\emptyset \vdash c : \ort{b}{\vnu = c} \ctinter \urt{b}{\vnu = c}$}
  \quad
  \mprooftr{54mm}
  {$\Gamma \ctcomma x{:}\tau_x \ctcomma f{:}(x{:}\ort{b}{\phi \land \vnu \prec x} \ctsarr \tau) \vdash e : \tau$}
  {T-Fix}
  {$\Gamma \vdash \zzfix{f}{x}{e} : x{:}\ort{b}{\phi} \ctsarr \tau$}
  \\[.8em]
  \mprooftr{34mm}
  {$\Code{Ty}(\primop) = \overline{x_i{:}\tau_i}\ctsarr \tau_y$ \quad
    $\forall i. \Gamma \vdash v_i : \tau_i$ }
  {T-AppOp}
  {$\Gamma \vdash \primop\;\overline{v_i} : \tau[\overline{x_i \mapsto v_i}]$}
  \quad
  \mprooftr{64mm}
  {$\forall i. \Gamma_i \vdash d_i(\overline{y}) : \ort{b}{\vnu = v}
    \interty \urt{b}{\vnu = v}$ \\
    $\forall i. \Gamma_i \vdash e_i : \tau_i$ }
  {T-Match}
  {$\ctbigoplus \overline{\Gamma_i} \vdash \match{v} \overline{d\ \overline{y} \to e}  : \ctbigoplus \overline{\tau_i}$}
  }
  \caption{Declarative typing rules for the core language.}
  \label{fig:declarative-typing-rules}
  \vspace*{-.1in}
\end{figure}

\paragraph{Variables and Constants} The typing rule for variables is standard;
a constant is allowed to have the most precise (singleton) type, which
has the same interpretation in both a demonic as well as angelic
setting.  In other words, the singleton type $\ort{b}{\vnu = c}
\ctinter \urt{b}{\vnu = c}$ coincides in both modalities because
there is only one possible value, so the safety and reachability
constraints are identical.


\paragraph{Bunched Typing} The typing rules are shown in
\autoref{fig:declarative-typing-rules}, where rules for
let-expressions, lambda functions, and applications have entangled and
independent versions, respectively.  Unlike standard bunched typing,
the entangled version does not require an explicit additive
conjunction in the type context, as shown below:
\begin{align*}
\mprooftr{40mm}
  {$\Gamma_1 \vdash e_1 : \tau_1 $ \quad $\Gamma_2 \ctcomma x{:}\tau_1 \vdash e_2 : \tau_2$}
  {T-Bunched-Let}
  {$\Gamma_1 \ctcomma \Gamma_2 \vdash \zlet{x}{e_1}{e_2} : \tau_2$}
\qquad
\mprooftr{26mm}
  {$\Gamma_1 \ctcomma \Gamma_2 \ctcomma \Gamma_2 \vdash e : \tau$}
  {T-Bunched-Contraction}
  {$\Gamma_1 \ctcomma \Gamma_2 \vdash e : \tau$}
\end{align*}\noindent
Unlike standard bunched typing, the entangled version does not require
an explicit additive conjunction in the type context. The
well-formedness condition requiring unique variable bindings means
that the substructural contraction rule
\textsc{T-Bunched-Contraction}, which duplicates bindings, is
incompatible with our system. Instead, our \textsc{T-Let} rule
directly entangles the newly added entangled context with the existing
context, without duplicating or consuming bindings.


\begin{example}[Let-expression] Consider the following example of typing for a let-expression:
\begin{align*}
    x{:}\ort{\Nat}{\vnu < 10} \vdash \zlet{y}{x - 3}{x - y} : \ort{\Nat}{\vnu \leq 3}
\end{align*}\noindent
Note that binding term $x - 3$ and body expression $x - y$ share the same variable $x$. If we apply standard bunched typing rule \textsc{T-Bunched-Let}, we require the following type context:
\begin{align*}
  \\&\underbrace{x{:}\ort{\Nat}{\vnu < 10}}_{\texttt{for $x - 3$}} \ctcomma \underbrace{x{:}\ort{\Nat}{\vnu < 10}}_{\texttt{for $x - y$}} \vdash \zlet{y}{x - 3}{x - y} : \ort{\Nat}{\vnu \leq 3}
\end{align*}\noindent
which requires using rule \textsc{T-Bunched-Contraction} to ``duplicate'' binding $x$. However, our \textsc{T-Let} rule optimizes this ``duplicate-then-consume'' pattern to yield the following derivation:
\begin{align*}
  &x{:}\ort{\Nat}{\vnu < 10} \vdash x - 3 : \ort{\Nat}{\vnu < 7}
  \\&\underline{x{:}\ort{\Nat}{\vnu < 10} \ctcomma y{:}\ort{\Nat}{\vnu < 7} \vdash x - y : \ort{\Nat}{\vnu \leq 3}}
  \\&x{:}\ort{\Nat}{\vnu < 10} \vdash \zlet{y}{x - 3}{x - y} : \ort{\Nat}{\vnu \leq 3}
\end{align*}\noindent
\end{example}

\paragraph{Subtyping}  Subsumption rules for both types (\textsc{T-Sub}) and
type contexts (\textsc{T-CtxSub}) use a \emph{semantic
  subtyping} relation. The type subtyping relation $\Gamma
\vdash\tau_1 \subty \tau_2$ is standard, i.e., inhabitants of the
denotation of $\tau_1$ are also inhabitants of the denotation of
$\tau_2$ under the denotation of context $\Gamma$. The context
subtyping relation $\Gamma_1 \ctsubseteq{X} \Gamma_2$ follows the
intuition in previous examples, where we restrict comparison to a set
of variables $X$.  For example,
\begin{align*}
  x{:}\ort{\Nat}{\vnu > 0} \ctcomma y{:}\urt{\Nat}{\top}
  \ctsubseteq{\{y\}} y{:}\urt{\Nat}{\top} \ctcomma z{:}\urt{\Nat}{\top}
\end{align*}
Rule \textsc{T-CtxSub} only requires the context-subtyping relation
to hold over free variables of the term and type. Further details are
given in~\autoref{sec:meta}.

\begin{example}[Substructural rules] We can derive the following substructural rules from \textsc{T-CtxSub}:
  \begin{align*}
    \mprooftr{25mm}
      {$\Gamma_2 \vdash e : \tau$}
      {T-Frame}
      {$\Gamma_1 \ctstar \Gamma_2 \vdash e : \tau$}
    \quad
    \mprooftr{26mm}
      {$\Gamma_2 \vdash e : \tau$}
      {T-Weakening}
      {$\Gamma_1 \ctcomma \Gamma_2 \vdash e : \tau$}
    \quad
    \mprooftr{25mm}
      {$\Gamma_2 \ctstar \Gamma_1 \vdash e : \tau$}
      {T-Exchange}
      {$\Gamma_1 \ctstar \Gamma_2 \vdash e : \tau$}
    \end{align*}\noindent
The variable-set constraint in context subtyping relation is important.
For example, framing rule \textsc{T-Frame} actually requires that $\Gamma_1 \ctstar \Gamma_2 \ctsubseteq{\Code{FV}(e) \cup \Code{FV}(\tau)} \Gamma_2$ holds, which is silently forced by $\Gamma_2 \vdash e : \tau$, requiring that the domain of $\Gamma_2$ is a subset of $\Code{FV}(e) \cup \Code{FV}(\tau)$.
\end{example}

\paragraph{Fixpoint Functions} As in standard refinement type systems,
we require all well-typed terms to be total. Thus, rule \textsc{T-Fix}
forces its first argument to have only base type and to always
decrease according to some well-founded relation $\prec$. This
decreasing constraint means the first parameter type cannot have
angelic modality; otherwise, the argument of recursive call
$v_{\Code{arg}}$ is required to have the following type:
\begin{align*}
  x{:}\urt{b}{\phi} \vdash v_{\Code{arg}} : \urt{b}{\phi} \interty
  \ort{b}{\vnu \prec x}
\end{align*}\noindent
where type $\urt{b}{\phi}$ requires $v_{\Code{arg}}$ to reach all
values that $x$ can reach; however, type $\ort{b}{\vnu \prec x}$
disallows $v_{\Code{arg}}$ from reaching the value of $x$. Such an
argument does not exist, which means recursive calls can never happen
when the parameter type has an angelic modality. Moreover, the
parameter type for a recursive function cannot be an independent arrow
type, i.e.,
\begin{align*}
  \Gamma \ctcomma x{:}\tau_x \ctcomma f{:}(x{:}\ort{b}{\phi \land \vnu \prec x} \ctwand \tau) \vdash e : \tau \quad \judgfail
\end{align*}\noindent
The parameter type of $f$ in a recursive call must reference $x$ via
the well-foundedness constraint $\nu \prec x$, which establishes a
dependency between the argument of the recursive call and the original
parameter $x$.  This dependency makes the independence requirement of
$\ctwand$ impossible to satisfy, since $\ctwand$ requires the argument
context to be independent of the closure context, whereas well-founded
induction inherently requires the opposite. Although \textsc{T-Fix}
therefore restricts the function type of recursive calls to use
$\ctsarr$, the subsumption rule \textsc{T-Sub} still allows a
recursive function to be typed against arbitrary function types after
the fact, provided the subtyping relation holds.

\begin{example}[Fixpoint Function] The following fixpoint function defines natural number subtraction discussed in \autoref{sec:intro}.
\begin{align*}
  \vdash &\zzfix{f}{x}{\zzlam{y}{\zif{x = 0}{0}{\zif{y = 0}{x}{f\;(x - 1)\;(y - 1)}}}}
  \\& : x{:}\ort{\Nat}{\top} \ctsarr y{:}\ort{\Nat}{\top} \ctsarr \ort{\Nat}{\vnu = x -_{\Nat} y}
\end{align*}\noindent
With a stronger function type that enables type checking with rule \textsc{T-Fix}, we can convert it into the type shown in \autoref{sec:intro} via  \textsc{T-Sub}.
\begin{align*}
  \vdash ... :  x{:}\urt{\Nat}{\vnu > 0} \ctwand y{:}\urt{\Nat}{\vnu > 0} \ctwand \urt{\Nat}{\vnu > 0}
\end{align*}\noindent
\end{example}

\paragraph{Datatypes and pattern matching} The last two rules
in \autoref{fig:declarative-typing-rules} type data constructors and
primitive operators. Rule \textsc{T-AppOp} is similar to the
application rule but relies on an auxiliary function $\Code{Ty}$ to
assign types to primitive operators and datatype constructors in the
core language. We require these primitive types to be precise as in
\textsc{T-Const}; for example:
\begin{align*}
  \Code{Ty}(\Code{O}) &= \ort{\Nat}{\vnu = 0} \interty \urt{\Nat}{\vnu = 0}
  \\\Code{Ty}(\Code{S}) &= n{:}\ort{\Nat}{\top} \ctsarr \ort{\Nat}{\vnu = n+1} \interty \urt{\Nat}{\vnu = n + 1}
  \\\Code{Ty}(+) &= x{:}\ort{\Nat}{\top} \ctsarr y{:}\ort{\Nat}{\top} \ctsarr \ort{\Nat}{\vnu = x+y} \interty \urt{\Nat}{\vnu = x + y}
\end{align*}\noindent
With these precise primitive types, rule \textsc{T-AppOp} can forward-infer constraints of operator application results. As discussed previously, rule \textsc{T-Match} requires type context to be divided into a sum of contexts that align with \emph{reachable} control flows, and then the sum of corresponding result types becomes the result type of the whole pattern-matching expression. The reachability condition is guaranteed by $\Gamma_i$ making datatype $d_i(\overline{y})$ have a precise type equal to the matched value $v$, which requires the types of variables $\overline{y}$ to have correct precise types for further type-checking of control-flow branches.

\begin{example}[Pattern matching] To illustrate \textsc{T-Match},
  consider how we can type \ref{ex:oplusWeak} from the previous
  section, after desugaring its conditional expression.
\begin{align*}
  x{:}\urt{\Nat}{\top} \vdash \zmatchNat{x}{3}{y}{5} : \ort{\Nat}{\vnu < 4} \ctoplus \ort{\Nat}{\top}
\end{align*}
Here, we want to show that the angelic context $x{:}\urt{\Nat}{\top}$
ensures that it is possible for \ref{ex:oplusWeak} to reach some value
less than $4$. To do so, we first partition the type context using
\textsc{T-CtxSub} %
{\small %
  \begin{align*}
    &x{:}\urt{\Nat}{\top} \ctsubseteq{\{x\}} \\
    & \quad
      \underbrace{x{:}\ort{\Nat}{\vnu = 0}\interty \urt{\Nat}{\vnu =
      0}}_{\Gamma_{\Code{trueBranch}}} \ctoplus
      \underbrace{x{:}\ort{\Nat}{\vnu > 0}\interty \urt{\Nat}{\vnu > 0}
      \ctcomma y{:}\ort{\Nat}{\vnu = x -
      1}}_{\Gamma_{\Code{falseBranch}}}
  \end{align*}}\noindent
into type contexts for each branch that align with the first
assumption of \textsc{T-Match}:
\begin{align*}
  &\Gamma_{\Code{trueBranch}} \vdash 0 : \ort{\Nat}{\vnu = 0} \interty \urt{\Nat}{\vnu = 0}
  \\&\Gamma_{\Code{falseBranch}} \vdash \Code{S}\;y : \ort{\Nat}{\vnu = x} \interty \urt{\Nat}{\vnu = x}
\end{align*}\noindent
We can then type each branch under these contexts to show that
\ref{ex:oplusWeak} has the target type:
\begin{align*}
  &\Gamma_{\Code{trueBranch}} \vdash 3 : \ort{\Nat}{\vnu < 4}
  \\&\underline{\Gamma_{\Code{falseBranch}} \vdash 5 : \ort{\Nat}{\top}}
  \\&\Gamma_{\Code{trueBranch}} \ctoplus \Gamma_{\Code{falseBranch}} \vdash  \zmatchNat{x}{3}{y}{5} : \ort{\Nat}{\vnu < 4} \ctoplus \ort{\Nat}{\top}
\end{align*}\noindent
On the other hand, we are unable to establish this under a weaker
typing context:
\begin{align*}
  x{:}\urt{\Nat}{\vnu > 0} \vdash \zmatchNat{x}{3}{y}{5} :
  \ort{\Nat}{\vnu < 4} \ctoplus \ort{\Nat}{\top} \quad \judgfail
\end{align*}\noindent
Here, $x{:}\urt{\Nat}{\vnu > 0}$ admits a world containing
environments that only reach the false branch, and thus cannot be
decomposed into two contexts that separately both branches.  This rule
also also enables the standard refinement type judgement with
unreachable control flows:
\begin{align*}
  &x{:}\ort{\Nat}{\vnu > 0} \ctcomma y{:}\ort{\Nat}{\vnu = x - 1}\vdash 5 : \ort{\Nat}{\vnu > 4}
  \\&x{:}\ort{\Nat}{\vnu > 0} \ctcomma y{:}\ort{\Nat}{\vnu = x - 1}\vdash \Code{S}\;y : \ort{\Nat}{\vnu = x} \interty \urt{\Nat}{\vnu = x}
  \\&\underline{x{:}\ort{\Nat}{\vnu > 0} \ctsubseteq{\{x\}} x{:}\ort{\Nat}{\vnu > 0} \ctcomma y{:}\ort{\Nat}{\vnu = x - 1}}
  \\&x{:}\ort{\Nat}{\vnu > 0} \vdash \zmatchNat{x}{3}{y}{5} : \ort{\Nat}{\vnu > 4} \quad \judgok
\end{align*}\noindent
where we adapt rule \textsc{T-Match} with a sum of a single
type. \BD{We need to clean this up a bit -- it is weird to leave off
  the true branch.}
\end{example}
